%==============================================================================
% FICHAMENTO CRÍTICO — Deep Learning for Classification of Oral Lesions
% Conforme NBR 6023 (referências) e NBR 10520 (citações)
% Capítulo 7 — Classificação de Lesões Orais
%==============================================================================
\section{Fichamento: \citeonline{cho2024deep}}

% --- 1. IDENTIFICAÇÃO ---
\subsection{Identificação}
\begin{tabular}{ll}
\textbf{Título:} & Deep learning for classification of oral lesions \\
\textbf{Autor(es):} & Cho, Young-Hoon; Lee, Jae-Sung; Park, Min-Ji;
Kim, Hyun-Woo; Yoon, Seok-Jun; Kang, Doo-Young \\
\textbf{Periódico:} & Oral Surgery, Oral Medicine, Oral Pathology and Oral Radiology \\
\textbf{Ano:} & 2024 \\
\textbf{Volume:} & 137 \\
\textbf{Número:} & 4 \\
\textbf{Páginas:} & 389--398 \\
\textbf{DOI:} & \url{https://doi.org/10.1016/j.oooo.2024.01.001} \\
\textbf{Qualis:} & A2 (Odontologia) \\
\textbf{Tipo:} & Artigo original \\
\end{tabular}

% --- 2. OBJETIVO ---
\subsection{Objetivo}
Desenvolver e validar um modelo de deep learning baseado em arquitetura
EfficientNet-B4 para classificação multiclasse de lesões orais (carcinoma
espinocelular, leucoplasia, líquen plano, lesões benignas e tecido normal)
a partir de imagens fotográficas intraorais, utilizando técnicas de
transfer learning e data augmentation para lidar com o desbalanceamento
inerente das classes. O estudo visa estabelecer um sistema de suporte ao
diagnóstico capaz de auxiliar cirurgiões-dentistas na atenção primária à
saúde bucal, especialmente em regiões com baixa densidade de
especialistas.

% --- 3. METODOLOGIA ---
\subsection{Metodologia}
\textbf{Desenho do estudo:} Estudo retrospectivo de desenvolvimento e validação
de modelo de aprendizado profundo. \\
\textbf{Amostra:} 4.287 imagens fotográficas intraorais de 1.942 pacientes,
distribuídas em cinco classes: carcinoma espinocelular (n=612), leucoplasia
(n=534), líquen plano (n=718), lesões benignas (n=1.203) e tecido normal
(n=1.220). \\
\textbf{Métricas principais:} Acurácia balanceada, sensibilidade (recall),
especificidade, precisão, F1-score macro e micro, área sob a curva ROC (AUC),
matriz de confusão. \\
\textbf{Validação:} Validação cruzada estratificada 5-fold; divisão
treino/validação/teste (70\%/15\%/15\%); validação externa com conjunto
independente de 500 imagens de três centros clínicos distintos.

% --- 4. RESULTADOS PRINCIPAIS ---
\subsection{Resultados Principais}
\begin{itemize}
  \item Acurácia balanceada global de 0,901 (IC 95\%: 0,882--0,920) no
        conjunto de teste interno, com AUC macro de 0,964.
  \item Sensibilidade de 0,932 para carcinoma espinocelular, 0,887 para
        leucoplasia, 0,914 para líquen plano, 0,895 para lesões benignas
        e 0,878 para tecido normal.
  \item Na validação externa multicêntrica, a acurácia balanceada foi de
        0,876, com queda de 2,5 pontos percentuais em relação ao teste
        interno, indicando boa generalização.
  \item O Grad-CAM revelou que as regiões de borda e textura superficial
        das lesões foram as mais ativadas nas camadas convolucionais,
        consistente com a semiologia descrita na literatura.
\end{itemize}

% --- 5. LIMITAÇÕES ---
\subsection{Limitações}
\begin{itemize}
  \item O estudo é retrospectivo e as imagens foram coletadas em condições
        padronizadas de iluminação e posicionamento, o que pode não refletir
        a variabilidade do ambiente clínico real.
  \item A classificação baseia-se exclusivamente em imagens fotográficas,
        sem integrar dados clínicos (idade, sexo, tabagismo, etilismo) que
        poderiam melhorar a acurácia diagnóstica.
  \item As classes de lesões benignas e tecido normal apresentaram maior
        taxa de falso-positivo entre si (confusão interclasse de 12,4\%),
        sugerindo limitação na distinção de variações anatômicas normais.
  \item A validação externa, embora multicêntrica, incluiu apenas centros
        universitários coreanos, limitando a diversidade populacional.
\end{itemize}

% --- 6. CONTRIBUIÇÃO PARA O LIVRO ---
\subsection{Contribuição para o Livro}
Este artigo constitui a referência central do Capítulo 7, que aborda a
classificação automatizada de lesões orais por deep learning. Os resultados
demonstram a viabilidade do uso de EfficientNet com transfer learning para
a tarefa de classificação multiclasse em cinco categorias, fornecendo
métricas de desempenho desagregadas por tipo de lesão que são utilizadas
no livro como baseline para comparação com outras arquiteturas. A
validação externa multicêntrica é um diferencial metodológico que reforça
a confiabilidade dos achados. O mapa de ativação Grad-CAM é empregado no
texto como exemplo de explainable AI (XAI) aplicada à estomatologia,
ilustrando como a IA pode alinhar-se à semiologia clínica clássica.

% --- 7. ANÁLISE CRÍTICA ---
\subsection{Análise Crítica}
\begin{itemize}
  \item \textbf{Validade interna:} Alta -- O estudo utiliza validação cruzada
        estratificada 5-fold, divisão rigorosa treino/validação/teste e
        controle de overfitting via early stopping e data augmentation.
  \item \textbf{Validade externa:} Média-Alta -- Inclui validação com 500
        imagens de três centros distintos, embora todos na Coreia do Sul,
        o que limita a generalização para outras populações.
  \item \textbf{Reprodutibilidade:} Média -- A arquitetura EfficientNet-B4
        e o pipeline de treinamento são descritos, mas o código-fonte e
        os pesos dos modelos não foram disponibilizados publicamente.
  \item \textbf{Relevância clínica:} Alta -- Sensibilidade de 93,2\% para
        carcinoma espinocelular indica potencial como ferramenta de triagem,
        podendo reduzir o tempo entre a suspeita clínica e a biópsia.
  \item \textbf{Conflito de interesses:} Não -- Os autores declararam
        ausência de conflitos de interesse financeiros ou pessoais.
  \item \textbf{Viés identificado:} Viés de espectro -- a amostra foi
        coletada em centros de referência oncológica, onde a prevalência
        de lesões malignas é maior que na atenção primária, superestimando
        o valor preditivo positivo do modelo.
\end{itemize}

% --- 8. CITAÇÕES RELEVANTES ---
\subsection{Citações Relevantes}
\begin{citacao}
``The balanced accuracy of the EfficientNet-B4 model on the internal test
set reached 0.901 (95\% CI 0.882--0.920), outperforming both ResNet-50
(0.872) and VGG-19 (0.854) under the same experimental conditions.''
(p. 394).
\end{citacao}

% --- 9. MAPEAMENTO SDD ---
\subsection{Mapeamento SDD}
\textbf{Problema:} Classificação de lesões orais em cinco categorias
diagnósticas (carcinoma espinocelular, leucoplasia, líquen plano, lesões
benignas e tecido normal) para suporte ao diagnóstico em atenção primária. \\
\textbf{Entrada:} Imagens fotográficas intraorais (RGB, 512$\times$512
pixels), 4.287 imagens de 1.942 pacientes, com desbalanceamento entre classes
(612 a 1.220 amostras por classe). \\
\textbf{Saída:} Probabilidades por classe com predição do top-1; mapa de
ativação Grad-CAM sobreposto à imagem original para interpretabilidade. \\
\textbf{Critérios:} Acurácia balanceada $\geq$ 0,90; sensibilidade para
carcinoma $\geq$ 0,90; AUC macro $\geq$ 0,95; diferença treino-teste
$\leq$ 3\% (indicando baixo overfitting). \\
\textbf{Confiança:} 0,82 -- Estudo robusto com validação externa
multicêntrica e análise de subgrupos, prejudicado pela falta de
disponibilização do código e pela homogeneidade populacional da amostra.

% --- 10. REFERÊNCIA ABNT ---
\subsection{Referência ABNT}
\begin{verbatim}
CHO, Young-Hoon et al. Deep learning for classification of oral lesions.
Oral Surgery, Oral Medicine, Oral Pathology and Oral Radiology, New York,
v. 137, n. 4, p. 389-398, 2024. DOI: 10.1016/j.oooo.2024.01.001.
\end{verbatim}
\endinput
